% Claim proof file for row claim_3_26 -- "AcyclicNonCollidersBlockable".
% Auto-generated stub by scaffold/solving_current_row.py at row start. A
% manager agent dedicated to writing the TeX proof fills in the body below.
% The proof must:
%   - Stay close to the lecture notes' proof if one exists (always search
%     the LN first for a `\begin{proof}` block immediately following the
%     claim; copy / lightly edit when present).
%   - Otherwise use the LN paradigm and cite prior claims by their `ref`.
%   - Be self-contained enough that a mathematician can follow it without
%     re-reading the LN.
%   - Have no `sorry`, `omitted`, or "obvious" shortcuts.
%
% Compiles standalone via the `subfiles` package.

\documentclass[main]{subfiles}

\begin{document}
% ref: claim_3_26 (proof, with statement restated for readability)
% title: AcyclicNonCollidersBlockable
\def\rowref{claim\_3\_26}
\def\rowtitle{AcyclicNonCollidersBlockable}
\phantomsection\label{claim_3_26}
%
% Cross-references: see claim_<ref>_statement_<title>.tex in this folder
% for the corresponding statement.

% --- statement (restated) ---------------------------------------------------
\begin{Rem}
    If a CDMG $G$ is acyclic then all non-colliders are blockable.
    So, the partial condition for $i\sigma$-separation
    ``a blockable non-collider in $C$''
    can be simplified to ``(any) non-collider in $C$''.
    
    So in the acyclic case we can simplify the notion of $i\sigma$-separation, 
    which we will refer to as \emph{$id$-separation}.
    However, in the non-acyclic setting $id$-separation (``(any) non-collider in $C$'') and $i\sigma$-separation (``a blockable non-collider in $C$'') are clearly not equivalent.

    It turned out that in the non-acyclic case $i\sigma$-separation is the more general concept (and as said above it also captures the acyclic case equivalently well), see \cite{FM17,FM18,FM20,BFPM21}.
    We will first focus on CADMGs (acyclic) for which we can restrict ourselves to the somewhat simpler $id$-separation.
    Later, we will pick up $i\sigma$-separation again when we deal with cycles.
\end{Rem}

% --- proof ------------------------------------------------------------------
\begin{proof}
% TODO: write the proof body.
\end{proof}

\end{document}
